\documentclass[a4paper,9pt]{article}

\usepackage{fullpage}
\usepackage[all]{xy}
\usepackage[utf8]{inputenc}
\usepackage{amsmath,amsthm}
\usepackage{amsfonts}
%\usepackage{algorithm}
\usepackage{graphicx}
%\usepackage[compatible,noend]{algpseudocode}
\usepackage{latexsym}
\usepackage{float,color}
\usepackage{multirow}
\usepackage{comment}

%%%%%%%%%%%%%%% COMMANDES %%%%%%%%%%%%%%%%%%%%%%%%%%%

\newcommand{\itemb}{\item[$\bullet$]}

%% Les ensembles de modelisation
\newcommand{\Ki}{\mathcal{K}}
\newcommand{\Ci}{\mathcal{C}}
\newcommand{\Mi}{\mathcal{M}}
\newcommand{\Ei}{\mathcal{E}}
\newcommand{\Di}{\mathcal{D}}

\newcommand{\Fd}{\mathbb{F}}
\newcommand{\Znat}{\mathbb{Z}}
\newcommand{\Nnat}{\mathbb{N}}

%%%%%%%%%%%%%%% ENVIRONEMENTS %%%%%%%%%%%%%%%%%%%%%%%


\theoremstyle{plain}
\newtheorem{lemme}{Lemme}
\newtheorem{algorithme}{ Algorithme }
\newtheorem{corolaire}{Corollaire}
\newtheorem{propriete}{Propri\'et\'e}
\newtheorem{proposition}{Proposition}
\newtheorem{theoreme} {Th\'eor\`eme}
\newtheorem{definition}{D\'efinition}
%\newtheorem*{correction}{Correction}
\excludecomment{correction}

\theoremstyle{definition}
\newtheorem*{notation}{Notation}
\newtheorem*{complexite}{Complexité}
\newtheorem{exercice}{Exercice}
%\theoremstyle{remark}
\newtheorem{remarque}{Remarque}
\newtheorem{exemple}{Exemple}

%0%%%%%%%%%%%% DOCUMENTS %%%%%%%%%%%%%%%%%%%%%%%%%%%%


\begin{document}



\hbox{\raisebox{0.4em}{\vrule depth 0pt height 0.5pt width \textwidth}}
\noindent \textbf{Université de Perpignan} \hfill  L3 Info \\
 \hfill C. Legrand \hfill  Année \the\year \\
\smallskip
\begin{center}
{
  {\scshape \large TD4 Sécurité }  \\
 Sécurité des réseaux 
}



\end{center}
\hbox{\raisebox{0.4em}{\vrule depth 0pt height 0.9pt width \textwidth}}
%\hline

\bigskip



\begin{exercice}[Vol de session TCP]
  Un attaquant Eve espionne une connexion \verb|telnet| entre Alice et Bob. Notez que dans l'application telnet lors d'un envoi d'un Alice$\to$Bob d'un message M, celui-ci est renvoyé en echo Bob$\to$Alice pour acquittement. 

 Eve forge un paquet TCP pour insérer la commande \verb|\n echo HACKED \n| dans le flux de données. Le dernier échange de paquets avant s'insertion est illlustré ci-dessous. Complétez la figure ci-dessous  avec le paquet inséré et les paquets qui suivront entre Alice et Bob.
  \begin{center}
\includegraphics[width=12cm]{vol-session.pdf}
    \end{center}
\end{exercice}


\begin{correction}
\includegraphics[width=11cm]{corrige-vol-session.pdf}
 
\end{correction}



\begin{exercice}[Trafic réseau suspect]
  Conrad Ministrateur, un informaticien qui travaille pour une banque suisse est responsable du bon fonctionnement et de la sécurité du réseau de son employeur. Un matin, lors d'un contrôle de routine, il remarque du trafic suspect à destination du serveur qui gère la base de données des comptes bancaires des clients.
  \bigskip
  
Voici un extrait de ce trafic:

\begin{verbatim}
  02:01:27 arp reply 192.168.94.1 is-at 0:3:47:48:5c:ee
  02:01:27 arp reply 192.168.94.1 is-at 0:3:47:48:5c:ee
  02:01:27 arp reply 192.168.94.1 is-at 0:3:47:48:5c:ee
  02:01:28 arp reply 192.168.94.1 is-at 0:3:47:48:5c:ee
  02:01:28 arp reply 192.168.94.1 is-at 0:3:47:48:5c:ee
  02:01:28 arp reply 192.168.94.1 is-at 0:3:47:48:5c:ee
  02:01:28 arp reply 192.168.94.1 is-at 0:3:47:48:5c:ee
  02:01:28 arp reply 192.168.94.1 is-at 0:3:47:48:5c:ee
  02:01:28 arp reply 192.168.94.1 is-at 0:3:47:48:5c:ee
  02:01:28 arp reply 192.168.94.1 is-at 0:3:47:48:5c:ee
  02:01:28 arp reply 192.168.94.1 is-at 0:3:47:48:5c:ee
  02:01:28 arp reply 192.168.94.1 is-at 0:3:47:48:5c:ee
  02:01:28 arp reply 192.168.94.1 is-at 0:3:47:48:5c:ee
  02:01:28 arp reply 192.168.94.1 is-at 0:3:47:48:5c:ee
  02:01:28 arp reply 192.168.94.1 is-at 0:3:47:48:5c:ee
  02:01:28 arp reply 192.168.94.1 is-at 0:3:47:48:5c:ee
  02:01:28 arp reply 192.168.94.1 is-at 0:3:47:48:5c:ee
  02:01:28 arp reply 192.168.94.1 is-at 0:3:47:48:5c:ee
  02:01:29 arp reply 192.168.94.1 is-at 0:3:47:48:5c:ee
\end{verbatim}

Sachant que l'adresse \verb|192.168.94.1| est celle de la passerelle qui offre un accès à Internet à ce serveur, interpréter ce trafic de manière concise. 
\end{exercice}



%\begin{exercice}
%On considère le réseau local suivant. Donnez les paquets que doits transmettre C  pour que les données envoyée par A soit transmise à C, et que les données provenant d'internet destinées à A soit transmise également transmise à C. 
%\end{exercice}


\begin{exercice}
On considère deux machines machine1 et machin2 dans un réseaux local connecté à internet par un routeur. On considère aussi un serveur web \verb|bienveillant.com| auquel veut accéder l'utilisateur de machine1 et \verb|malveillant.com|.  Les adresses IP et adresses MAC sont données ci-dessous:
\begin{center}
\includegraphics[width=7cm]{reseau-empoisonement.pdf}
  \end{center}

\begin{enumerate}
\item On suppose que la machine \verb|machine1| fait un ping vers l'adresse IP 192.168.1.2: il émet le paquet suivant  sur le réseau
\begin{verbatim}
 192.168.1.1 envoie [ARB who-as? 192.168.1.2] à FF:FF:FF:FF:FF:FF
\end{verbatim}

Expliquez comment l'attaquant sur machine2 doit faire pour que les données transmises par machine1 au routeur lui soit d'abord transmise. 
\item Expliquez ensuite comment l'attaque sur machine2 doit s'y prendre pour que machine1 envoie ses requetes web à \verb|site-malveillant.com| au lieu de \verb|mon-gentil-site.com|. 
\end{enumerate}

\end{exercice}


\begin{exercice}
  Jean-Michel Pyrhatt a remarqué qu'une requête de transfert de zone DNS (utilisé notamment pour fournir à un serveur DNS esclave l'information  concernant la ou les zones qu'il doit dessert) fait toujours 27 octets de long et que la réponse du serveur DNS \verb|dns.machinchose.ch| fait 745 octets. On admettra que les communications se font au moyen d'UDP.

  \begin{enumerate}
  \item Comment Jean-Michel Pyrhatt peut-il réaliser une attaque par déni de service contr le serveur \verb|pauvre.victime.com| ?
    \item Si Jean-Michel Pyrhatt possède une bande passante de 1Mbps et qu'il est capable de l'utiliser à 100\% pour son attaque, quelle sera la bande passante appartenant à \verb|pauvre.victime.com| qui sera utilisée par l'attaque ?

   \end{enumerate}
\end{exercice}


\begin{correction}
\begin{enumerate}
\item L'attaque consiste à envoyer des requêtes DNS \verb|dns.machinchose.ch| en usurpant l'adresse IP source du serveur \verb|pauvre.victime.com| à rendre inopérant. Le flux de message d'erreurs sera 30 fois plus important et inonderant le serveur cible.  

\item Si Jean-Michel Pyratt a une bande passante de 1Mbs, le serveur DNS génèrera $\sim 30$Mps qui occuperont la bande passante de \verb|pauvre.victime.com|.
\end{enumerate}
\end{correction}




\begin{exercice}[Attaque de Kevin Mitnick]
  Le trafic décrit ci-dessous a été observé dans un réseau (cas réel). De quoi s'agit-il ? Les paquets TCP sont décrits à l'aide de lignes ayant le format suivant :
\begin{verbatim}
  14.18.37.26 alice.513 > bob.514: P 1382727010.1382727011(2) ack 2024384001
\end{verbatim}
On a :
\begin{itemize}
\item l'heure \verb|14.18.37.26|
\item l'adresse source (\verb|alice|)
\item le port source (\verb|513|)
\item  l'adresse de detination (\verb|bob|)
\item le port de destination (\verb|514|)
\item un drapeau éventuel (\verb|S|=SYN, \verb|P|=PUSH, \verb|F|=FIN, \verb|R|=RESET, \verb|.| = pas de drapeau.
\item Les numéro de séquence du premier octet et du dernier octet du datagramme (1382727010.1382727011).
\item le nombre d'octets transmis \verb|2|.
\item Un éventuel acquittement des octets reçus (\verb|ACK|).
  \item Le numéro de séquence du prochain octet attendu (\verb|2024384001|).
\end{itemize}


\begin{verbatim}
14:18:22.516699 mallory.600 > alice.513: S 1382726960:1382726960(0) win 4096
14:18:22.566069 mallory.601 > alice.513: S 1382726961:1382726961(0) win 4096
14:18:22.744477 mallory.602 > alice.513: S 1382726962:1382726962(0) win 4096
14:18:22.830111 mallory.603 > alice.513: S 1382726963:1382726963(0) win 4096
14:18:22.886128 mallory.604 > alice.513: S 1382726964:1382726964(0) win 4096
14:18:22.943514 mallory.605 > alice.513: S 1382726965:1382726965(0) win 4096
14:18:23.002715 mallory.606 > alice.513: S 1382726966:1382726966(0) win 4096
14:18:23.103275 mallory.607 > alice.513: S 1382726967:1382726967(0) win 4096
14:18:23.162781 mallory.608 > alice.513: S 1382726968:1382726968(0) win 4096
14:18:23.225384 mallory.609 > alice.513: S 1382726969:1382726969(0) win 4096
14:18:23.282625 mallory.610 > alice.513: S 1382726970:1382726970(0) win 4096
14:18:23.342657 mallory.611 > alice.513: S 1382726971:1382726971(0) win 4096
14:18:23.403083 mallory.612 > alice.513: S 1382726972:1382726972(0) win 4096
14:18:23.903700 mallory.613 > alice.513: S 1382726973:1382726973(0) win 4096
14:18:24.003252 mallory.614 > alice.513: S 1382726974:1382726974(0) win 4096
14:18:24.084827 mallory.615 > alice.513: S 1382726975:1382726975(0) win 4096
14:18:24.142774 mallory.616 > alice.513: S 1382726976:1382726976(0) win 4096
14:18:24.203195 mallory.617 > alice.513: S 1382726977:1382726977(0) win 4096
14:18:24.294773 mallory.618 > alice.513: S 1382726978:1382726978(0) win 4096
14:18:24.382841 mallory.619 > alice.513: S 1382726979:1382726979(0) win 4096
14:18:24.443309 mallory.620 > alice.513: S 1382726980:1382726980(0) win 4096
14:18:24.643249 mallory.621 > alice.513: S 1382726981:1382726981(0) win 4096
14:18:24.906546 mallory.622 > alice.513: S 1382726982:1382726982(0) win 4096
14:18:24.963768 mallory.623 > alice.513: S 1382726983:1382726983(0) win 4096
14:18:25.022853 mallory.624 > alice.513: S 1382726984:1382726984(0) win 4096
14:18:25.153536 mallory.625 > alice.513: S 1382726985:1382726985(0) win 4096
14:18:25.400869 mallory.626 > alice.513: S 1382726986:1382726986(0) win 4096
14:18:25.483127 mallory.627 > alice.513: S 1382726987:1382726987(0) win 4096
14:18:25.599582 mallory.628 > alice.513: S 1382726988:1382726988(0) win 4096
14:18:25.653131 mallory.629 > alice.513: S 1382726989:1382726989(0) win 4096
......
4:18:25.906002 kevin.1000 > bob.514: S 1382726990:1382726990(0) win 4096
14:18:26.094731 bob.514 > kevin.1000: S 2021824000:2021824000(0) ack 1382726991 win 4096
14:18:26.172394 kevin.1000 > bob.514: R 1382726991:1382726991(0) win 0
14:18:26.507560 kevin.999 > bob.514: S 1382726991:1382726991(0) win 4096
14:18:26.694691 bob.514 > kevin.999: S 2021952000:2021952000(0) ack 1382726992 win 4096
14:18:26.775037 kevin.999 > bob.514: R 1382726992:1382726992(0) win 0
14:18:26.775395 kevin.999 > bob.514: R 1382726992:1382726992(0) win 0
14:18:27.014050 kevin.998 > bob.514: S 1382726992:1382726992(0) win 4096
14:18:27.174846 bob.514 > kevin.998: S 2022080000:2022080000(0) ack 1382726993 win 4096
14:18:27.251840 kevin.998 > bob.514: R 1382726993:1382726993(0) win 0
14:18:27.544069 kevin.997 > bob.514: S 1382726993:1382726993(0) win 4096
14:18:27.714932 bob.514 > kevin.997: S 2022208000:2022208000(0) ack 1382726994 win 4096
14:18:27.794456 kevin.997 > bob.514: R 1382726994:1382726994(0) win 0
14:18:28.054114 kevin.996 > bob.514: S 1382726994:1382726994(0) win 4096
14:18:28.224935 bob.514 > kevin.996: S 2022336000:2022336000(0) ack 1382726995 win 4096
14:18:28.305578 kevin.996 > bob.514: R 1382726995:1382726995(0) win 0
14:18:28.564333 kevin.995 > bob.514: S 1382726995:1382726995(0) win 4096
14:18:28.734953 bob.514 > kevin.995: S 2022464000:2022464000(0) ack 1382726996 win 4096
14:18:28.811591 kevin.995 > bob.514: R 1382726996:1382726996(0) win 0
14:18:29.074990 kevin.994 > bob.514: S 1382726996:1382726996(0) win 4096
14:18:29.274572 bob.514 > kevin.994: S 2022592000:2022592000(0) ack 1382726997 win 4096
14:18:29.354139 kevin.994 > bob.514: R 1382726997:1382726997(0) win 0
14:18:29.354616 kevin.994 > bob.514: R 1382726997:1382726997(0) win 0
14:18:29.584705 kevin.993 > bob.514: S 1382726997:1382726997(0) win 4096
14:18:29.755054 bob.514 > kevin.993: S 2022720000:2022720000(0) ack 1382726998 win 4096
14:18:29.840372 kevin.993 > bob.514: R 1382726998:1382726998(0) win 0
14:18:30.094299 kevin.992 > bob.514: S 1382726998:1382726998(0) win 4096
14:18:30.265684 bob.514 > kevin.992: S 2022848000:2022848000(0) ack 1382726999 win 4096
14:18:30.342506 kevin.992 > bob.514: R 1382726999:1382726999(0) win 0
14:18:30.604547 kevin.991 > bob.514: S 1382726999:1382726999(0) win 4096
14:18:30.775232 bob.514 > kevin.991: S 2022976000:2022976000(0) ack 1382727000 win 4096
14:18:30.852084 kevin.991 > bob.514: R 1382727000:1382727000(0) win 0
14:18:31.115036 kevin.990 > bob.514: S 1382727000:1382727000(0) win 4096
14:18:31.284694 bob.514 > kevin.990: S 2023104000:2023104000(0) ack 1382727001 win 4096
14:18:31.361684 kevin.990 > bob.514: R 1382727001:1382727001(0) win 0
14:18:31.627817 kevin.989 > bob.514: S 1382727001:1382727001(0) win 4096
14:18:31.795260 bob.514 > kevin.989: S 2023232000:2023232000(0) ack 1382727002 win 4096
14:18:31.873056 kevin.989 > bob.514: R 1382727002:1382727002(0) win 0
14:18:32.164597 kevin.988 > bob.514: S 1382727002:1382727002(0) win 4096
14:18:32.335373 bob.514 > kevin.988: S 2023360000:2023360000(0) ack 1382727003 win 4096
14:18:32.413041 kevin.988 > bob.514: R 1382727003:1382727003(0) win 0
14:18:32.674779 kevin.987 > bob.514: S 1382727003:1382727003(0) win 4096
14:18:32.845373 bob.514 > kevin.987: S 2023488000:2023488000(0) ack 1382727004 win 4096
14:18:32.922158 kevin.987 > bob.514: R 1382727004:1382727004(0) win 0
14:18:33.184839 kevin.986 > bob.514: S 1382727004:1382727004(0) win 4096
14:18:33.355505 bob.514 > kevin.986: S 2023616000:2023616000(0) ack 1382727005 win 4096
14:18:33.435221 kevin.986 > bob.514: R 1382727005:1382727005(0) win 0
14:18:33.695170 kevin.985 > bob.514: S 1382727005:1382727005(0) win 4096
14:18:33.985966 bob.514 > kevin.985: S 2023744000:2023744000(0) ack 1382727006 win 4096
14:18:34.062407 kevin.985 > bob.514: R 1382727006:1382727006(0) win 0
14:18:34.204953 kevin.984 > bob.514: S 1382727006:1382727006(0) win 4096
14:18:34.375641 bob.514 > kevin.984: S 2023872000:2023872000(0) ack 1382727007 win 4096
14:18:34.452830 kevin.984 > bob.514: R 1382727007:1382727007(0) win 0
14:18:34.714996 kevin.983 > bob.514: S 1382727007:1382727007(0) win 4096
14:18:34.885071 bob.514 > kevin.983: S 2024000000:2024000000(0) ack 1382727008 win 4096
14:18:34.962030 kevin.983 > bob.514: R 1382727008:1382727008(0) win 0
14:18:35.225869 kevin.982 > bob.514: S 1382727008:1382727008(0) win 4096
14:18:35.395723 bob.514 > kevin.982: S 2024128000:2024128000(0) ack 1382727009 win 4096
14:18:35.472150 kevin.982 > bob.514: R 1382727009:1382727009(0) win 0
14:18:35.735077 kevin.981 > bob.514: S 1382727009:1382727009(0) win 4096
14:18:35.905684 bob.514 > kevin.981: S 2024256000:2024256000(0) ack 1382727010 win 4096
14:18:35.983078 kevin.981 > bob.514: R 1382727010:1382727010(0) win 0
....
14:18:36.245045 alice.513 > bob.514: S 1382727010:1382727010(0) win 4096
14:18:36.755522 alice.513 > bob.514: .                          ack 2024384001 win 4096
14:18:37.265404 alice.513 > bob.514: P 1382727010:1382727011(2) ack 2024384001 win 4096
14:18:37.775872 alice.513 > bob.514: P 1382727012:1382727016(5) ack 2024384001 win 4096
14:18:38.287404 alice.513 > bob.514: P 1382727017:1382727041(25) ack 2024384001 win 4096
14:18:36.755522 alice.513 > bob.514: .                          ack 2024384002 win 4096
14:18:36.755522 alice.513 > bob.514: .                          ack 2024384003 win 4096
14:18:43.165874 alice.513 > bob.514: F 1382727041:1382727041(0) ack 2024384003 win 4096
14:18:52.179922 alice.513 > bob.514: R 1382727043:1382727043(0) win 4096
14:18:52.236452 alice.513 > bob.514: R 1382727044:1382727044(0) win 4096
.....
14:18:52.298431 mallory.600 > alice.513: R 1382726960:1382726960(0) win 4096
14:18:52.363877 mallory.601 > alice.513: R 1382726961:1382726961(0) win 4096
14:18:52.416916 mallory.602 > alice.513: R 1382726962:1382726962(0) win 4096
14:18:52.476873 mallory.603 > alice.513: R 1382726963:1382726963(0) win 4096
14:18:52.536573 mallory.604 > alice.513: R 1382726964:1382726964(0) win 4096
14:18:52.600899 mallory.605 > alice.513: R 1382726965:1382726965(0) win 4096
14:18:52.660231 mallory.606 > alice.513: R 1382726966:1382726966(0) win 4096
14:18:52.717495 mallory.607 > alice.513: R 1382726967:1382726967(0) win 4096
14:18:52.776502 mallory.608 > alice.513: R 1382726968:1382726968(0) win 4096
14:18:52.836536 mallory.609 > alice.513: R 1382726969:1382726969(0) win 4096
14:18:52.937317 mallory.610 > alice.513: R 1382726970:1382726970(0) win 4096
14:18:52.996777 mallory.611 > alice.513: R 1382726971:1382726971(0) win 4096
14:18:53.056758 mallory.612 > alice.513: R 1382726972:1382726972(0) win 4096
14:18:53.116850 mallory.613 > alice.513: R 1382726973:1382726973(0) win 4096
14:18:53.177515 mallory.614 > alice.513: R 1382726974:1382726974(0) win 4096
14:18:53.238496 mallory.615 > alice.513: R 1382726975:1382726975(0) win 4096
14:18:53.297163 mallory.616 > alice.513: R 1382726976:1382726976(0) win 4096
14:18:53.365988 mallory.617 > alice.513: R 1382726977:1382726977(0) win 4096
14:18:53.437287 mallory.618 > alice.513: R 1382726978:1382726978(0) win 4096
14:18:53.496789 mallory.619 > alice.513: R 1382726979:1382726979(0) win 4096
14:18:53.556753 mallory.620 > alice.513: R 1382726980:1382726980(0) win 4096
14:18:53.616954 mallory.621 > alice.513: R 1382726981:1382726981(0) win 4096
14:18:53.676828 mallory.622 > alice.513: R 1382726982:1382726982(0) win 4096
14:18:53.736734 mallory.623 > alice.513: R 1382726983:1382726983(0) win 4096
14:18:53.796732 mallory.624 > alice.513: R 1382726984:1382726984(0) win 4096
14:18:53.867543 mallory.625 > alice.513: R 1382726985:1382726985(0) win 4096
14:18:53.917466 mallory.626 > alice.513: R 1382726986:1382726986(0) win 4096
14:18:53.976769 mallory.627 > alice.513: R 1382726987:1382726987(0) win 4096
14:18:54.039039 mallory.628 > alice.513: R 1382726988:1382726988(0) win 4096
14:18:54.097093 mallory.629 > alice.513: R 1382726989:1382726989(0) win 4096
\end{verbatim}

\end{exercice}

\begin{correction}
\begin{itemize}
\item La première séquence de messages mallory$\to$alice est une attaque par déni de service de type SYN flooding. Le but est de rendre muette alice qui ne pourra pas répondre à bob. 

\item La séquence de messages kevin$\to$bob est une séquence d'ouverture de connexion, directement d'une annulation (Reset). Elle a pour but à l'attaquant kevin d'idendifier comment les numéros ISN sont choisis lors de l'ouverture d'une connexion.

\item La séquence de messages alice$\to$bob sont des messages envoyés par kevin en usurpant l'identité d'alice (l'attaque se fait en aveugle). Il se sert de la connaissance de l'ISN pour pouvoir injecter des faux messages : 
\begin{itemize}
\item kevin prédit l'ISN de bob et  répond (en aveugle) au SYN-ACK,
\item kevin injecte des données 25 octets (commande malveillante reçu comme venant d'alice machine de confiance) 
\item kevin termine la communication avec un FIN et des acquittements.
\end{itemize}
\item La dernière séquence de messages mallory$\to$alice sont des message Reset, qui annule les demandes d'ouvertures du SYN flooding afin de permettre à alice de communiquer à nouveau. 
\end{itemize}
\end{correction}



\begin{exercice}
On considère un pare-feu sans état qui abrite la machine 203.167.75.1. L'utilisateur de cette machine souhaite naviguer sur le Web, recevoir des connexions \verb|telnet| provenant de l'extérieur du réseau interne, ainsi que recevoir et initier des connexions SSH de et vers internet. Les servurs HTTP, \verb|telnet| et SSH utilisent repsectivement les ports 80,23,22.

\'Ecrire la table de filtrage du pare-feu en utilisant le critère de filtrage sur les paquets SYN (flag SYN à 1 et flag ACK à 0).
\end{exercice}

\begin{correction}
Voici la table 
\begin{center}
\begin{tabular}{|c|c|c|c|c|c|c|}
\hline
regle & @ src & Port source &   @ dest.& Port dest. & SYN & action \\
\hline
\hline
1 & * & * & 203.167.75.1 & 22 & * & autorisé  \\
2 & * & 22 & 203.167.75.1 & * & non & autorisé  \\
3 & * & * & 203.167.75.1 & 23 & * & autorisé \\
4 & * & 80 & 203.167.75.1 & * & non & autorisé  \\
\hline
5 &  203.167.75.1 &  22 & *  & * & non & autorisé  \\
6 & 203.167.75.1 &  *  & *  & 22 & * & autorisé  \\
7 & 203.167.75.1 &  23 & *  & * & non & autorisé  \\
8 & 203.167.75.1 &  * & *  & 80 & * & autorisé \\
\hline
* & *& *& *& * & bloqué \\
\hline
\end{tabular}
\end{center}

\begin{itemize}
\item Regle 1: autorise la connexion (SYN) et requete (ACK) de machine exterieur au serveur SSH de la machine 203.167.75.1.
\item Regle 2: autorise les réponses uniquement (SYN-ACK,ACK) de serveur SSH extérieur à la machine 203.167.75.1.
\item Regle 3: autorise la connexion (SYN) et requete (ACK) d'une machine extérieure sur le serveur telnet de  203.167.75.1.
\item Regle 4: autorise la réponse (SYN-ACK,ACK) d'un serveur web à la machine 203.167.75.1.
\item Regle 5: autorise les réponses (SYN-ACK,ACK) du serveur SSH de 203.167.75.1 à une machine extérieur.
\item Regle 6: autorise la connexion et réponse/requete d'un client ssh  203.167.75.1 vers un serveur ssh d'une machine extérieur.
\item Regle 7: autorise les réponses du serveur telnet sur 203.167.75.1  à une machine de l'extérieur. 
\item Regle 8:  Autorise les connexion et requete d'un client web sur  203.167.75.1 vers un  serveur web externe.
\end{itemize}

\end{correction}




\begin{exercice}
On considère la règle de parefeu avec état  pour le réseau donné ci-dessous. 
La règle A du pare-feu permet aux machines du LAN privé d'accéder à la DMZ 2, alors que la règle C devait l'interdire. Comment remédier à cela?
\begin{center}
\includegraphics[width=10cm]{LAN-securise}
\end{center}

\begin{center}
\begin{tabular}{|c|c|c|c|c|c|c|}
\hline
Règle & @ src & @ dest. & Protocole & Port source & Port dest. & Action \\
\hline
\hline
A & Toutes & DMZ 2 & TCP & Tous & 80 & Autorisé \\
\hline
B & LAN & DMZ 1 & TCP & Tous & 25 & Autorisé \\
\hline
C & LAN & Toutes & TCP & Tous & Tous & Refusé \\
\hline
E & Tous & Tous & Tous & Tous & Tous & Refusé \\
\hline
\end{tabular}
\end{center}
\end{exercice}

%\begin{comment}
\begin{correction}
Une solution consiste à passer la règle A en troisième position:

\begin{center}
\begin{tabular}{|c|c|c|c|c|c|c|}
\hline
Règle & @ src & @ dest. & Protocole & Port source & Port dest. & Action \\
\hline
\hline
B & LAN & DMZ 1 & TCP & Tous & 25 & Autorisé \\
\hline
C & LAN & Toutes & TCP & Tous & Tous & Refusé \\
\hline
A & Toutes & DMZ 2 & TCP & Tous & 80 & Autorisé \\
\hline
E & Tous & Tous & Tous & Tous & Tous & Refusé \\
\hline
\end{tabular}
\end{center}
La règle B permet aux machines du LAN d'accéder au DMZ 1. La règle C interdit tout autre trafic en provenance du LAN. La règle A n'a plus d'influence sur le trafic du LAN et permet aux machines externes d'accéder à DMZ 2.

\end{correction}
%\end{comment}

\begin{exercice}
  On considère l'architecture réseau donnée ci-dessous. Les Firewall sont des pare-feu avec état.  Le serveur web a pour adresse IP 10.0.0.2 et les proxy SMTP, HTTP, DNS les adresses IP respectires 10.0.0.25, 10.0.0.80, 10.0.0.53. Le serveur web doit être accessible par le réseau interne. On désignera dmz-proxy les adresses IP de tous les proxy. Donnez les règles du Firewall1 qui permette d'assurer ces fonctionnalités.
  \begin{center}
\includegraphics[width=8cm]{dmz-sandwitch.pdf}
    \end{center}
\end{exercice}

%\end{document}


\begin{correction}

  $\;$
  
  
   \begin{center}
    \begin{tabular}{|c|c|c|c|c|c|c|}
      \hline
      Règle & source & port & destination & port & protocole & action \\
      \hline
      1 & * & * & 10.0.0.25 & 25 & tcp & autorisé \\
      2 & * & * & 10.0.0.80 & 80 & tcp & autorisé \\
      3 & * & * & 10.0.0.53 & 53 & udp & autorisé \\
      4 & * & * & dmz-proxy & * & * & refusé \\
      \hline
      5 &  10.0.0.25 & * & * & 25 & tcp & autorisé \\
      6 &  10.0.0.80 & * & * & 80 & tcp & autorisé \\
      7 &  10.0.0.53 & * & * & 53 & udp & autorisé \\
      8 & dmz-proxy & * & * & * & * & refusé \\      
      \hline
      9 & * & * & 10.0.0.2 & 80 & tcp & autorisé \\
      10 & * & * & dmz-web & * & * &  refusé \\
      \hline    
      11 & dmz-web & * & *  &  * & * & refusé \\
      \hline
      12 & * & * & * & * & * & refusé \\
      \hline
    \end{tabular}
  \end{center}

\end{correction}


\end{document}

